\section{boot.asm 中的 PSE 临时映射}

在第~\ref{ch:vmm}~章中我们已经知道，\texttt{boot.asm} 使用 4\,MiB PSE 大页
建立临时映射。这里我们详细分析为什么需要\textbf{同时}建立 identity mapping 和
high-half mapping。

\codefile{src/boot/boot.asm}
\begin{lstlisting}[style=nasmstyle, caption={boot.asm：构建双重 PSE 映射}]
; PSE PDE: [31:22]=phys_base | PS(bit7) | RW(bit1) | P(bit0)
;   flags = 0x83
xor ecx, ecx              ; ecx = 0..255
.fill_pse:
    mov eax, ecx
    shl eax, 22            ; eax = ecx * 4 MiB
    or  eax, 0x83          ; PS | RW | Present
    mov [boot_pd + ecx*4],         eax  ; identity: PD[i]
    mov [boot_pd + 768*4 + ecx*4], eax  ; high-half: PD[768+i]
    inc ecx
    cmp ecx, 256
    jb  .fill_pse
\end{lstlisting}

这段代码会产生如下映射：

\begin{center}
\begin{tabular}{lll}
  \toprule
  \textbf{PDE 范围} & \textbf{虚拟地址} & \textbf{映射到物理地址} \\
  \midrule
  PD[0..255]     & \hex{00000000}--\hex{3FFFFFFF} & \hex{00000000}--\hex{3FFFFFFF} \\
  PD[768..1023]  & \hex{C0000000}--\hex{FFFFFFFF} & \hex{00000000}--\hex{3FFFFFFF} \\
  \bottomrule
\end{tabular}
\end{center}

\begin{whybox}
为什么不能只建 high-half mapping，跳过 identity mapping？

因为在开启分页的那一瞬间，\texttt{EIP} 仍然指向物理地址 $\sim$\hex{001xxxxx}
（\texttt{boot.asm} 的 \texttt{.boot} 段链接在低地址）。
如果没有 identity mapping，\reg{CR0.PG} 置 1 后 CPU 取下一条指令时会
Page Fault → Double Fault → Triple Fault → 重启。

\textbf{时序回顾}：
\begin{enumerate}
  \item \texttt{mov cr0, eax} — 开启分页，此时 EIP $\approx$ \hex{001xxxxx}
  \item CPU 取下一条指令 → 虚拟地址 \hex{001xxxxx} → identity mapping 翻译 → 物理 \hex{001xxxxx} ✓
  \item \texttt{jmp \_start\_high} — 跳转到虚拟地址 \hex{C01xxxxx}
  \item CPU 取指 → \hex{C01xxxxx} → high-half mapping 翻译 → 物理 \hex{001xxxxx} ✓
  \item 从此以后，代码运行在 high-half，identity mapping 变成了"遗产"
\end{enumerate}
\end{whybox}

开启分页的关键指令序列：

\codefile{src/boot/boot.asm}
\begin{lstlisting}[style=nasmstyle, caption={boot.asm：开启分页并跳转到 high-half}]
; 加载页目录物理地址到 CR3
mov eax, boot_pd
mov cr3, eax

; 开启分页 (CR0.PG = 1)
; [CPU STATE] 从下一条指令开始，所有内存访问经过 MMU
mov eax, cr0
or  eax, 0x80000000     ; bit 31 = PG
mov cr0, eax

; --- 分页已开启，identity mapping 保证 EIP 仍然有效 ---

; 跳转到 high-half 虚拟地址
mov eax, _start_high    ; _start_high 链接在 0xC01xxxxx
jmp eax
\end{lstlisting}

\begin{cpustate}
执行 \texttt{mov cr0, eax}（PG=1）后：
\begin{itemize}
  \item CR3 → \texttt{boot\_pd}（物理地址）
  \item EIP 仍在低地址（$\sim$\hex{001xxxxx}），通过 PD[0] 的 identity mapping 继续运行
  \item \texttt{jmp \_start\_high} 后，EIP 切换到 \hex{C01xxxxx}，通过 PD[768] 的 high-half mapping 运行
  \item 此后内核栈、全局变量等均位于 \hex{C0xxxxxx}+ 虚拟地址
\end{itemize}
\end{cpustate}

% ============================================================================

